
\section{Background}\label{sec:pca} 
We commence our exposition by recalling several well-known algebraic structures: pointed convex algebras, categories enriched over them and the category of sub-Markov kernels which is the archetypical example of a category with such an enrichment.

\subsection{Pointed Convex Algebras}\label{ssec:pca} A \emph{pointed convex algebra} (pca) \cite{stone1949postulates} consists of a set $X$, a designated element $\star \in X$ and, for all $p$ in the open real interval $(0,1)$, a function $+_p \colon X\times X \to X$ such that, by fixing $\tilde{p}\defeq pq$ and $\tilde{q}\defeq \frac{p(1-q)}{1-pq}$, the following laws hold for all $x_1,x_2,x_3\in X$. % the following laws hold for .
\begin{equation}\label{eq:pca}(x_1+_q x_2)+_p x_3 =  x_1 +_{\tilde{p}} (x_2+_{\tilde{q}} x_3) \qquad x_1+_px_2=x_2+_{1-p}x_1 \qquad  x_1+_p x_1 = x_1\end{equation}
A morphism of pcas is a function preserving $\star$ and $+_p$. %Pointed convex algebras and their morphisms form a category, hereafter denoted by $\Cat{PCA}$.
We denote by $\Cat{PCA}$ the category of pcas and their morphisms.  

In any pca, $+_p$ can be defined for all $p\in[0,1]$ by setting $x+_1 y \defeq x$ and $x+_0 y \defeq y$. For all $n\in \mathbb{N}$, $p_1, \dots, p_n \in [0,1]$ such that $\sum_i p_i\leq 1$, one defines $\sum_{i=1}^n p_i\cdot (-)_i \colon X^n \to X$ inductively as
%\begin{equation}\label{eq:def somma n pca}
%    \sum_{i=1}^0 p_i\cdot (-)_i \defeq \star \qquad \sum_{i=1}^{n+1} p_i\cdot (-)_i \defeq (-)_{n+1} +_{p_{n+1}} \sum_{i=1}^n \frac{p_i}{1-p_{n+1}}\cdot (-)_i\text{.}
%\end{equation}
%
\begin{equation}\label{eq:def somma n pca}
    \sum_{i=1}^0 p_i\cdot (-)_i \defeq \star \qquad \sum_{i=1}^{n+1} p_i\cdot (-)_i \defeq (-)_{1} +_{p_{1}} \sum_{j=1}^n q_j \cdot (-)_j\text{.}
\end{equation}
where $(-)_j=(-)_{i+1}$ and $q_j$ is $\frac{p_{i+1}}{1-p_1}$ if $p_1 \neq 1$ and $0$ otherwise.
Note that when $n=1$, $\sum_{i=1}^1p_i \cdot (-)_i$ is, by definition, $(-)_1 +_{p_1} \star$. Since this operation will play a crucial role, we name it \emph{multiplication by a scalar} $p\in [0,1]$ and we fix the following notation. \[p\cdot x \defeq x+_p \star\] 
%These $n$-ary sums enjoy the following properties: (a) $\sum_{i=1}^n p_i\cdot x_i =x_j$ if $p_j =1$ and (b) $\sum_{i=1}^n p_1 \cdot ( \sum_{j=1}^n q_{i,j} \cdot x_j) = \sum_{j=1}^m(\sum_{i=1}^n p_i q_{i,j} x_i)$. Actually, one can prove (see \cite[Proposition 7]{bonchi2017power}) that $n$-ary sums with the properties (a) and (b) provide an alternative but equivalent presentation for pcas.

The standard example of a pca is provided by the set $\Dis(X)$ of finitely supported probability subdistributions over some set $X$. Recall that a finitely supported subdistribution over  $X$ is a function $d\colon X \to [0,1]$ such that $\sum_{x\in X}d(x)\leq1$ and $d(x)\neq 0$ for finitely many $x$. The distinguished element $\star$ in $\Dis(X)$ is the null subdistribution, i.e., $\star(x)\defeq 0$ for all $x\in X$; given $d_1,d_2\in \Dis(X)$, for all $p\in(0,1)$, $(d_1+_pd_2)(x) \defeq p\cdot d_1(x)+ (1-p)\cdot d_2(x)$. At this point, it is convenient to fix some extra notation: $\DisF(X)$ is the set of all finitely supported distributions (i.e., $\sum_{x\in X}d(x)=1$) and, for all $x\in X$, $\delta_x$ is the Dirac distribution at $x$ (i.e., $\delta_x(x')=1$ if $x=x'$ and $0$ otherwise).

Interestingly, the pointed convex algebra $\Dis(X)$  enjoys an additional property that will be relevant for our development: cancellativity. A pca $(X,+_p,\star)$ is \emph{cancellative} (or, in the terminology of \cite{sokolova2018termination}, cancellative at $\star$) if the following implication holds for all $x,y\in X$ and $p\in(0,1)$.
\begin{equation*}
p\cdot x = p\cdot y \Rightarrow x=y
\end{equation*}
\begin{remark}\label{rmk:collapsed}
The reader may have noticed that the axioms of a pca resemble the familiar laws of
associativity, commutativity, and idempotency for monoids. A crucial
difference, however, is that the distinguished element $\star$ is \emph{not} required
to act as a unit. Indeed, any pca that additionally satisfies the axiom
\begin{equation}\label{eq:unitstar}
x+_p \star = x
\end{equation}
collapses to a commutative idempotent monoid. More precisely, if one enriches
\eqref{eq:pca} with the condition above, one can easily show that
$
x+_p y = x+_q y$
for all $x,y\in X$ and $p,q\in(0,1)$,
so that the parameter $p$ becomes irrelevant. See, for instance, \cite{DBLP:journals/lmcs/BonchiSV22} for a detailed derivation.
\end{remark}



%\begin{remark}
%The reader may have noticed that the axioms of pca in  resemble associativity, commutativity and idempotency in monoids. The fact that $\star$ is \emph{not} forced to act as a unit is crucial. Indeed, those pcas that additionally satisfy the axiom
%\begin{equation}
%x+_p \star =x
%\end{equation}
%degenerate to commutative idempotent monoids: by adding to \eqref{eq:pca} the axiom above, one can easily prove that $x+_p y =x+_q y$ for all $x,y\in X$ and $p,q\in(0,1)$. See, e.g., \cite{DBLP:journals/lmcs/BonchiSV22} for a derivation.
%\end{remark}



\subsection{Categories enriched over pcas} A category $\Cat{C}$ is \emph{$\Cat{PCA}$-enriched} if every homset $\Cat{C}[X,Y]$ carries a pca structure and composition of arrows is a pcas morphism, namely that the following equalities hold for all $p\in(0,1)$ and properly typed arrows $e,f,g,h$.
\begin{equation}\label{eq:enr}e; (f+_p g) = (e;f)+_p (e;g) \qquad (f+_pg) ;h= (f;h +_p g;h) \qquad f;\star =\star= \star;f\end{equation}
 It is now convenient to note a few properties of scalar multiplication in $\Cat{PCA}$-enriched categories.
\begin{lemma}\label{lemma:initialproperties2}
In a $\Cat{PCA}$-enriched category the following properties hold:
\begin{enumerate}
\item $(p\cdot f) ; g =p\cdot (f;g)$ and $f;(p\cdot g) = p\cdot ( f;g)$ ; \label{lemma:p;}
\item $p\cdot (q \cdot f) =pq \cdot f$; \label{lemma:pq}
%\item $p \cdot \diagq{X} = \langle \id{X},\id{X}\rangle_{pq, p(1-q)} $.\label{lemma:pdiag}
%\item $\diagp{X}; (\id{X} \oplus q\cdot \id{X})=\langle \id{X},\id{X}\rangle_{p,(1-p)q}$ \label{lemma:diagp}
%tolto i due commentati sotto perché non si usano...
%\item $(\sum_{i=1}^{n}p_i\cdot f_i);h=\sum_{i=1}^{n}p_i\cdot (f_i;h) $ and $k;\sum_{i=1}^{n}p_i\cdot f_i= \sum_{i=1}^{n}p_i\cdot (h;f_i)$; \label{lemma sum ; h}
\item $q\cdot\sum_{i=1}^{n}p_i\cdot f_i= \sum_{i=1}^{n} (qp_i)\cdot f_i$;\label{lemma q per somma}
%\item $\sum_{i=1}^{n}{p_i\cdot(f_i+_p g_i)}= (\sum_{i=1}^{n}p_i\cdot f_i)+_p (\sum_{i=1}^{m}p_i\cdot g_i)$ \label{lemma:somma di +_p}
\end{enumerate}
\end{lemma}


A functor $F$ is $\Cat{PCA}$-enriched if it preserves the pca structure of each homset. $\Cat{PCA}$-enriched categories and functors form a category, denoted by $\Cat{PCACat}$.
All categories considered hereafter are tacitly assumed to be locally small: $\Cat{Cat}$ stands  for the category of locally small categories. Note that there is a forgetful functor $U\colon \Cat{PCACat} \to \Cat{Cat}$.

Any category $\Cat{C}$ gives rise to the $\Cat{PCA}$-enriched category $\Cat{C}^+$: objects of  $\Cat{C}^+$ are those of $\Cat{C}$; for all objects $X,Y$, the homset is defined as $\Cat{C}^+[X,Y] \defeq \Dis(\Cat{C}[X,Y])$. For $d_1\colon X \to Y$ and $d_2\colon Y \to Z$, their composition $d_1;d_2 \colon X\to Z$ is defined for all $h\in \Cat{C}[X,Z]$ as $d_1;d_2(h) \defeq \sum_{\{(f,g)\mid f;g=h\}}d_1(f) \cdot d_2(g)$; The identity $\id{X}\colon X \to X$ is given by $\delta_{id_X}$. One can easily see that $\Cat{C}^+$ is a $\Cat{PCA}$-enriched category. Moreover, the assignment $\Cat{C} \mapsto \Cat{C}^+$ gives rise to a functor $(-)^+\colon \Cat{Cat} \to \Cat{PCACat}$ which is left adjoint to the forgetful functor $U$. % (see e.g. \cite[Prop. 6.4.7]{borceux2} or \cite[Cor. 1]{villoria2024enriching}). % for a general result.

\begin{theorem}[\cite{borceux2,villoria2024enriching}]\label{thm:freeenriched}
 $(-)^+\colon \Cat{Cat} \to \Cat{PCACat}$ is left adjoint to $U\colon \Cat{PCACat} \to \Cat{Cat}$.
\end{theorem}
\begin{proof}%[Proof of Theorem~\ref{thm:freeenriched}]%\marginpar{To be moved in the appendix}\marginpar{Questo risultato di sicuro esiste in letteratura}
The result follows from more general results about enriching on arbitrary algebraic theories: see for instance \cite[Prop. 6.4.7]{borceux2} or \cite[Cor. 1]{villoria2024enriching}. For later use, it is convenient anyway to briefly illustrate the involved structures. We first show the definition of $F^+\colon \Cat{C}^+ \to \Cat{D}^+$ for a functor $F\colon \Cat{C} \to \Cat{D}$. For all objects $X$, $F^+$ is defined as $F^+(X)\defeq F(X)$; For all $d\in \Cat{C}^+[X,Y]$, $F^+(d)\in \Cat{D}^+[FX, FY]$ is defined for all $g\colon FX \to FY$ as $F^+(d)(g)\defeq \sum_{\{f\in \Cat{C}[X,Y]| F(f)=g\}}d(f)$.  One can easily check that the functor is $\Cat{PCA}$-enriched.


For all categories $\Cat{C}$, the unit $\eta_{\Cat{C}}\colon \Cat{C} \to \Cat{C}^+$ is the functor acting as identity on objects and mapping any arrow $f\in \Cat{C}[X,Y]$ into $\delta_f \in \Cat{C^+}[X,Y]$.

Given a $\Cat{PCA}$-enriched category $\Cat{D}$ and a functor $F\colon \Cat{C} \to U(\Cat{D})$, one can define a $\Cat{PCA}$-enriched functor $F^\sharp \colon \Cat{C}^+ \to \Cat{D}$ as follows: for all objects $X$, $F^\sharp(X)\defeq F(X)$, for all arrows $d\in \Cat{C}^+[X,Y]$, $F^\sharp(d)\defeq \sum_{f\in \Cat{C}[X,Y]}d(f)\cdot {F(f)}$. One can easily check that $\eta_{\Cat{C}} ; F^\sharp=F$.   

\end{proof}

\begin{example}\label{ex:PCA}
Let $\Cat{1}$ be a category with a single object and a single arrow. The $\Cat{PCA}$-enriched category $\Cat{1}^+$ has a single object, arrows are $p\in[0,1]$ and composition is given by multiplication.

For a set $A$, the set of words  over $A$ (hereafter denoted by $A^*$) carries the structure of a category with a single object.  The $\Cat{PCA}$-enriched category $(A^*)^+$ has a single object and arrows are subdistributions over $A^*$. The composition of $d_1,d_2\in \Dis(A^*)$ is defined for all words $w\in A^*$ as $d_1;d_2(w)\defeq \sum_{u;v=w}d_1(u)\cdot d_2(v)$. 
\end{example}

\subsection{The category of sub-Markov kernels}
Our main example of a $\Cat{PCA}$-enriched category is $\KlD$, the Kleisli category of the monad $\subdistr\colon \Sets \to \Sets$ (see e.g.,~\cite{hasuo2007generic}).
Objects are sets; morphisms \(f \colon X \to Y\) are functions \(X \to \subdistr(Y)\). % mapping each \(x \in X\) into a finitely supported subdistribution over $Y$.
  We often write \(f(y \mid x)\) for \(f(x)(y)\), as this number represents the probability that \(f\) returns \(y\) given the input \(x\).
  Identities \(\id{X} \colon X \to \subdistr(X)\) map each element \(x \in X\) to \(\delta_{x}\).
  For two functions \(f \colon X \to \subdistr(Y)\) and \(g \colon Y \to \subdistr(Z)\), their composition in $\KlD$ is defined as \(f ; g (z \mid x) \defeq \sum_{y \in Y} f(y \mid x) \cdot g(z \mid y)\).

\begin{proposition}
$\KlD$ is $\Cat{PCA}$-enriched.
\end{proposition}
\begin{proof}
While this result is well known, it is anyway convenient to illustrate its proof. For all objects $X,Y$ of $\KlD$, the hom-set $\KlD[X,Y]$ is a pointed convex algebra: for all $f,g\colon X \to Y$, $x\in X$, $y\in Y$ and $p\in (0,1)$,
\[f +_p g(y|x)\defeq p\cdot f(y|x) + (1-p)\cdot g(y|x) \qquad \star_{X,Y}(y|x) \defeq 0.\] Note that $\star_{X,Y}$ is the arrow which sends every $x$ into the null subdistribution over $Y$. An easy calculation shows that arrow composition is a morphism of pointed convex algebras, i.e., it preserves the operations $+_p$ and $\star_{X,Y}$. For instance $(f+_pg) ;h= (f;h +_p g;h)$ for all $f,g\colon X \to Y$ and $h\colon Y \to Z$ is obtained as follows:
\begin{align}
    ((f+_pg) ;h)(z|x) &= \sum_{y\in Y} (f+_pg)(y|x) \cdot h(z|y) \tag*{} \\
    &= \sum_{y\in Y} (p\cdot f(y|x) + (1-p)\cdot g(y|x)) \cdot h(z|y) \tag*{}\\
    &= p\cdot \sum_{y\in Y} f(y|x) \cdot h(z|y) + (1-p)\cdot \sum_{y\in Y} g(y|x) \cdot h(z|y) \tag*{}\\
    &= p\cdot (f;h)(z|x) + (1-p)\cdot (g;h)(z|x) \tag*{}\\
    &= (f;h +_p g;h)(z|x)\text{.}\tag*{}
    \end{align}
\end{proof}




In our work, it is fundamental that $\KlD$ carries two symmetric monoidal structures: $ (\KlD, \per, \uno)$ and $(\KlD, \piu, \zero)$. The monoidal product $\per$ is defined on objects as the cartesian product of sets, often denoted by $\times$, with unit the singleton $\uno \defeq \{\bullet\}$; $\piu$ is the disjoint union of sets with unit the empty set $\zero \defeq \{\}$. Hereafter we denote the disjoint union of two sets $X$ and $Y$ by $X\oplus Y \defeq \{(x,0) \mid x\in X\} \cup \{(y,1) \mid y \in Y\}$ where $0$ and $1$ are tags used to distinguish the set of provenance. For arrows  \(f \colon X \to Y\) and \(g \colon X' \to Y'\), $f\per g\colon X\otimes X' \to Y \otimes Y'$ and $f \piu g \colon X\oplus X' \to Y\oplus Y'$ are defined, for all $x\in X$, $x'\in X'$, $y\in Y$, $y'\in Y'$, $u\in X\oplus X'$ and $v\in Y\oplus Y'$, as follows.
\begin{equation}\label{ex:products}
\begin{aligned}
f \per g(y,y' \mid x,x') &\defeq f(y \mid x) \cdot g(y' \mid x')\\
f\piu g(v \mid u) &\defeq 
\begin{cases} 
f(y \mid x) & \text{if } u=(x,0) \text{ and } v=(y,0)\\ 
g(y' \mid x') & \text{if } u=(x',1) \text{ and } v=(y',1) \\  
0 & \text{otherwise} 
\end{cases}
\end{aligned}
\end{equation}
Symmetries $\symmt{X}{Y}\colon X \otimes Y \to Y \otimes X$ and $\symmp{X}{Y}\colon X \oplus Y \to Y \oplus X$ are defined as in $\Sets$. 

Actually, $ (\KlD, \piu, \per, \zero, \uno)$ forms a rig (also known as bimonoidal) category in the sense of \cite{laplaza_coherence_1972}. We will come back to rig categories in Section \ref{sec:tapediagrams}. Until then, we will focus only on the monoidal structure provided by $\oplus$. 
%


